\documentclass[10pt,a4paper]{amsart}
\usepackage{amsmath,amssymb,amsthm,graphicx}
% Retained minimal preamble for the Stacks Project chapter varieties.tex.
% Verbatim excerpt of the chapter preamble: Xy-pic setup (native preamble lines 14-18), the native theorem
% declarations with the mathematical macros (native preamble lines 154-217), and the project own no-op
% definitions of the reference/slogan/history display environments (native wrapper.tex lines 1-3) that the
% counted statement uses for its slogan. Copied mechanically; nothing added, removed or redefined.
% For commutative diagrams we use Xy-pic
\usepackage[all]{xy}

% We use 2cell for 2-commutative diagrams.
\xyoption{2cell}
% Theorem environments.
%
\theoremstyle{plain}
\newtheorem{theorem}[subsection]{Theorem}
\newtheorem{proposition}[subsection]{Proposition}
\newtheorem{lemma}[subsection]{Lemma}

\theoremstyle{definition}
\newtheorem{definition}[subsection]{Definition}
\newtheorem{example}[subsection]{Example}
\newtheorem{exercise}[subsection]{Exercise}
\newtheorem{situation}[subsection]{Situation}

\theoremstyle{remark}
\newtheorem{remark}[subsection]{Remark}
\newtheorem{remarks}[subsection]{Remarks}

\numberwithin{equation}{subsection}

% Macros
%
\def\lim{\mathop{\mathrm{lim}}\nolimits}
\def\colim{\mathop{\mathrm{colim}}\nolimits}
\def\Spec{\mathop{\mathrm{Spec}}}
\def\Hom{\mathop{\mathrm{Hom}}\nolimits}
\def\Ext{\mathop{\mathrm{Ext}}\nolimits}
\def\SheafHom{\mathop{\mathcal{H}\!\mathit{om}}\nolimits}
\def\SheafExt{\mathop{\mathcal{E}\!\mathit{xt}}\nolimits}
\def\Sch{\mathit{Sch}}
\def\Mor{\mathop{\mathrm{Mor}}\nolimits}
\def\Ob{\mathop{\mathrm{Ob}}\nolimits}
\def\Sh{\mathop{\mathit{Sh}}\nolimits}
\def\NL{\mathop{N\!L}\nolimits}
\def\CH{\mathop{\mathrm{CH}}\nolimits}
\def\proetale{{pro\text{-}\acute{e}tale}}
\def\etale{{\acute{e}tale}}
\def\QCoh{\mathit{QCoh}}
\def\Ker{\mathop{\mathrm{Ker}}}
\def\Im{\mathop{\mathrm{Im}}}
\def\Coker{\mathop{\mathrm{Coker}}}
\def\Coim{\mathop{\mathrm{Coim}}}

% Boxtimes
%
\DeclareMathSymbol{\boxtimes}{\mathbin}{AMSa}{"02}

%
% Macros for moduli stacks/spaces
%
\def\QCohstack{\mathcal{QC}\!\mathit{oh}}
\def\Cohstack{\mathcal{C}\!\mathit{oh}}
\def\Spacesstack{\mathcal{S}\!\mathit{paces}}
\def\Quotfunctor{\mathrm{Quot}}
\def\Hilbfunctor{\mathrm{Hilb}}
\def\Curvesstack{\mathcal{C}\!\mathit{urves}}
\def\Polarizedstack{\mathcal{P}\!\mathit{olarized}}
\def\Complexesstack{\mathcal{C}\!\mathit{omplexes}}
% \Pic is the operator that assigns to X its picard group, usage \Pic(X)
% \Picardstack_{X/B} denotes the Picard stack of X over B
% \Picardfunctor_{X/B} denotes the Picard functor of X over B
\def\Pic{\mathop{\mathrm{Pic}}\nolimits}
\def\Picardstack{\mathcal{P}\!\mathit{ic}}
\def\Picardfunctor{\mathrm{Pic}}
\def\Deformationcategory{\mathcal{D}\!\mathit{ef}}
\newenvironment{reference}{}{}
\newenvironment{slogan}{}{}
\newenvironment{history}{}{}
\usepackage{cleveref}
\crefname{originallemma}{original lemma}{original lemmas}
\Crefname{originallemma}{Original lemma}{Original lemmas}
\makeatletter
\newlabel{algebra-lemma-when-surjective-local}{{0E8M}{source}{}{sourcepointer-1}{}}
\newlabel{algebra-lemma-when-surjective-local@cref}{{[originallemma][0][]0E8M}{[1][0][]source}{}{}{}}
\makeatother
\crefname{originallemma}{original lemma}{original lemmas}
\Crefname{originallemma}{Original lemma}{Original lemmas}
\makeatletter
\newlabel{coherent-lemma-proper-finite-fibre-finite-in-neighbourhood}{{02OH}{source}{}{sourcepointer-2}{}}
\newlabel{coherent-lemma-proper-finite-fibre-finite-in-neighbourhood@cref}{{[originallemma][0][]02OH}{[1][0][]source}{}{}{}}
\makeatother
\crefname{originaldefinition}{original definition}{original definitions}
\Crefname{originaldefinition}{Original definition}{Original definitions}
\makeatletter
\newlabel{constructions-example-projective-space}{{0FCY}{source}{}{sourcepointer-3}{}}
\newlabel{constructions-example-projective-space@cref}{{[originaldefinition][0][]0FCY}{[1][0][]source}{}{}{}}
\makeatother
\crefname{originallemma}{original lemma}{original lemmas}
\Crefname{originallemma}{Original lemma}{Original lemmas}
\makeatletter
\newlabel{lemma-tangent-space-rational-point}{{0B2E}{source}{}{sourcepointer-4}{}}
\newlabel{lemma-tangent-space-rational-point@cref}{{[originallemma][0][]0B2E}{[1][0][]source}{}{}{}}
\makeatother
\crefname{originallemma}{original lemma}{original lemmas}
\Crefname{originallemma}{Original lemma}{Original lemmas}
\makeatletter
\newlabel{morphisms-lemma-image-proper-scheme-closed}{{01W6}{source}{}{sourcepointer-5}{}}
\newlabel{morphisms-lemma-image-proper-scheme-closed@cref}{{[originallemma][0][]01W6}{[1][0][]source}{}{}{}}
\makeatother
\title{Very ampleness from separating closed points and nonzero tangent vectors by sections}
\author{The Stacks Project Authors}
\date{stacks-project:varieties.tex:local-snapshot-20261011; GNU FDL 1.2 or later}
\begin{document}
\maketitle
\noindent\textit{Editorial note.} All delivered mathematical text is the author's own, copied from the pinned Stacks Project chapter file varieties.tex; no statement or proof step was rewritten, reordered or removed.\par\medskip
\noindent\textit{Editorial note.} The proof contains two bounded asides, '(computation omitted)' at the injectivity on closed points and '(small detail omitted)' at the singleton fibre, both reproduced verbatim; the surrounding steps are otherwise complete as written.\par\medskip
\noindent\textit{Editorial note.} The counted spans carry 5 navigation stubs and no bibliography citation.\par\medskip
\noindent\textit{Editorial note.} The preamble is the chapter's retained minimal preamble (sources/stacks-varieties/retained-preamble.tex lines 1-77), a verbatim excerpt of the native preamble.\par\medskip
\noindent\textit{Editorial note.} The label lemma-separate-points-tangent-vectors is defined only in varieties.tex, so no cross-chapter restatement had to be weighed for this card.\par\medskip
\noindent\textit{Allowed standard prerequisites (editorial).}\begin{itemize}
\item Canonical morphism to projective space attached to a subspace of global sections. Constructions, Example constructions-example-projective-space, invoked at proof lines 4328-4335. Supplies the morphism phi\_\{L,V\}: X to P(V) whose closed-immersion property is at stake.
\item Image of a proper scheme is closed and a proper morphism with finite fibres is finite. Morphisms, Lemma morphisms-lemma-image-proper-scheme-closed and Cohomology of Schemes, Lemma coherent-lemma-proper-finite-fibre-finite-in-neighbourhood, invoked at proof lines 4336-4342. Supplies properness of phi and finiteness once the fibres are singletons.
\item Tangent space of a projective space at a rational point and the cotangent description of surjectivity. Same chapter, Lemma lemma-tangent-space-rational-point, invoked at proof lines 4374-4377. Supplies that injectivity of the tangent map is surjectivity of m\_p/m\_p\textasciicircum{}2 to m\_x/m\_x\textasciicircum{}2.
\item Local criterion for surjectivity of a map of local rings with surjective cotangent map. Algebra, Lemma algebra-lemma-when-surjective-local, invoked at proof lines 4379-4381. Supplies the final surjectivity of the stalk map from the finite module structure and the surjection on cotangent spaces.
\end{itemize}\medskip
\ifdefined\hypertarget\hypertarget{sourcepointer-1}{}\fi
\noindent\textit{Source reference.} Original lemma 0E8M: sources/stacks-chow/tags.txt, source lines 17822--17822. Pointer outside the counted span, invoked once for context; the load-bearing argument of the counted statement is the author's proof reproduced in full in the proof span.\par\medskip
\ifdefined\hypertarget\hypertarget{sourcepointer-2}{}\fi
\noindent\textit{Source reference.} Original lemma 02OH: sources/stacks-chow/tags.txt, source lines 3516--3516. Pointer outside the counted span, invoked once for context; the load-bearing argument of the counted statement is the author's proof reproduced in full in the proof span.\par\medskip
\ifdefined\hypertarget\hypertarget{sourcepointer-3}{}\fi
\noindent\textit{Source reference.} Original definition 0FCY: sources/stacks-chow/tags.txt, source lines 19198--19198. Pointer outside the counted span, invoked once for context; the load-bearing argument of the counted statement is the author's proof reproduced in full in the proof span.\par\medskip
\ifdefined\hypertarget\hypertarget{sourcepointer-4}{}\fi
\noindent\textit{Source reference.} Original lemma 0B2E: sources/stacks-varieties/pinned-original-source.tex, source lines 3059--3059. Pointer outside the counted span, invoked once for context; the load-bearing argument of the counted statement is the author's proof reproduced in full in the proof span.\par\medskip
\ifdefined\hypertarget\hypertarget{sourcepointer-5}{}\fi
\noindent\textit{Source reference.} Original lemma 01W6: sources/stacks-chow/tags.txt, source lines 2471--2471. Pointer outside the counted span, invoked once for context; the load-bearing argument of the counted statement is the author's proof reproduced in full in the proof span.\par\medskip
\begin{lemma}
\label{lemma-separate-points-tangent-vectors}
Let $k$ be an algebraically closed field.
Let $X$ be a proper $k$-scheme.
Let $\mathcal{L}$ be an invertible $\mathcal{O}_X$-module.
Let $V \subset H^0(X, \mathcal{L})$ be a $k$-subvector space. If
\begin{enumerate}
\item for every pair of distinct closed points $x, y \in X$
there is a section $s \in V$ which vanishes at $x$ but not at $y$, and
\item for every closed point $x \in X$ and nonzero tangent vector
$\theta \in T_{X/k, x}$ there exists a section $s \in V$
which vanishes at $x$ but whose pullback by $\theta$ is nonzero,
\end{enumerate}
then $\mathcal{L}$ is very ample and the canonical morphism
$\varphi_{\mathcal{L}, V} : X \to \mathbf{P}(V)$
is a closed immersion.
\end{lemma}

\begin{proof}
Condition (1) implies in particular that the elements of $V$ generate
$\mathcal{L}$ over $X$. Hence we get a canonical morphism
$$
\varphi = \varphi_{\mathcal{L}, V} :
X
\longrightarrow
\mathbf{P}(V)
$$
by Constructions, Example \ref{constructions-example-projective-space}.
The morphism $\varphi$ is proper by
Morphisms, Lemma \ref{morphisms-lemma-image-proper-scheme-closed}.
By (1) the map $\varphi$ is injective on closed points
(computation omitted).
In particular, the fibre over any closed point of $\mathbf{P}(V)$
is a singleton (small detail omitted). Thus we see that
$\varphi$ is finite, for example use Cohomology of Schemes, Lemma
\ref{coherent-lemma-proper-finite-fibre-finite-in-neighbourhood}.
To finish the proof it suffices to show that the map
$$
\varphi^\sharp :
\mathcal{O}_{\mathbf{P}(V)}
\longrightarrow
\varphi_*\mathcal{O}_X
$$
is surjective. This we may check on stalks at closed points.
Let $x \in X$ be a closed point with image the closed
point $p = \varphi(x) \in \mathbf{P}(V)$. Since
$\varphi^{-1}(\{p\}) = \{x\}$ by (1) and since $\varphi$
is proper (hence closed), we see that $\varphi^{-1}(U)$
runs through a fundamental system of open neighbourhoods
of $x$ as $U$ runs through a fundamental system of open neighbourhoods of $p$.
We conclude that on stalks at $p$ we obtain the map
$$
\varphi^\sharp_x :
\mathcal{O}_{\mathbf{P}(V), p}
\longrightarrow
\mathcal{O}_{X, x}
$$
In particular, $\mathcal{O}_{X, x}$ is a finite
$\mathcal{O}_{\mathbf{P}(V), p}$-module.
Moreover, the residue fields of $x$ and $p$ are equal to $k$
(as $k$ is algebraically closed -- use the Hilbert Nullstellensatz).
Finally, condition (2) implies that the map
$$
T_{X/k, x} \longrightarrow T_{\mathbf{P}(V)/k, p}
$$
is injective since any nonzero $\theta$ in the kernel of this map
couldn't possibly satisfy the conclusion of (2).
In terms of the map of local rings above this means that
$$
\mathfrak m_p/\mathfrak m_p^2 \longrightarrow
\mathfrak m_x/\mathfrak m_x^2
$$
is surjective, see Lemma \ref{lemma-tangent-space-rational-point}.
Now the proof is finished by applying
Algebra, Lemma \ref{algebra-lemma-when-surjective-local}.
\end{proof}

\end{document}
